% Source: https://artofproblemsolving.com/community/c3h2479201_flipping_inequality_signs
% Stable identity: aops-topic-2479201-q01
\textbf{(a)} We have $b^2-a^2=(b-a)(b+a)$ with $b-a>0$. If $0\le a<b$, then $b+a>0$, so $b^2-a^2>0$. If $a<b\le0$, then $b+a<0$, so $b^2-a^2<0$, i.e. $a^2>b^2$. (So squaring keeps the sign when both sides are nonnegative and flips it when both are nonpositive.)

\textbf{(b)} $-1<2$ with $1<4$; $\ -2<2$ with $4=4$; $\ -3<1$ with $9>1$. So when the sides have opposite signs, squaring gives no information.

\textbf{(c)} Here $b>0$, so $\sqrt a+\sqrt b>0$ and
\[\sqrt b-\sqrt a=\dfrac{b-a}{\sqrt b+\sqrt a}>0.\]
(So taking square roots of nonnegative numbers never flips the sign.)

\textbf{(d)} The left side requires $x\ge-2$.
\begin{itemize}
\item If $-2\le x<0$: the left side is $\ge0>x$, so the inequality holds.
\item If $x\ge0$: both sides are nonnegative, so by (a) and (c) the inequality is equivalent to $x+2>x^2$, i.e. $(x-2)(x+1)<0$, i.e. $-1<x<2$. Together with $x\ge0$ this gives $0\le x<2$.
\end{itemize}
So the solution set is $\boxed{-2\le x<2}$. (Squaring without the case split would wrongly lose $-2\le x\le-1$.)
